\documentclass{article}

% Pacotes
\usepackage[utf8]{inputenc}
\usepackage[brazil]{babel}
\usepackage[letterpaper,top=2.5cm,bottom=2.5cm,left=3cm,right=3cm,marginparwidth=1.75cm]{geometry}
\usepackage{amsmath, graphicx, subcaption}
\usepackage[colorlinks=true, allcolors=blue]{hyperref}

\usepackage{listings}
\usepackage{xcolor}
\usepackage{tabularx}
\usepackage{array}
\usepackage{amssymb}
\usepackage{extarrows}
\usepackage{url}
\usepackage{hyperref}
\usepackage{booktabs}
\usepackage{float}
\usepackage{graphicx}
\usepackage{enumitem}


% Definições de estilo para Python
\lstdefinestyle{pythonstyle}{
      language=Python,
      basicstyle=\ttfamily\small,
      keywordstyle=\color{blue}\bfseries,
      commentstyle=\color{gray}\itshape,
      stringstyle=\color{orange},
      numbers=left,
      numberstyle=\tiny,
      stepnumber=1,
      numbersep=10pt,
      backgroundcolor=\color{white},
      frame=single,
      breaklines=true,
      tabsize=4,
      showstringspaces=false,
}


\begin{document}

% Capa personalizada
\begin{titlepage}
      \centering
      \includegraphics[width=4cm]{images/fig_logo_cefetmg.pdf}\par
      \vspace{0.5cm}

      {\LARGE Centro Federal de Educação Tecnológica de Minas Gerais \par}
      \vspace{0.2cm}

      {\Large Engenharia de Computação \\ Computação Gráfica \par}
      \vspace{4cm}

      {\Huge \textbf{Relatório de Desenvolvimento} \par}
      \vspace{0.5cm}

      {\LARGE Renderização da Bandeira Brasileira em Espaço RGB \par}
      \vspace{3cm}

      {\LARGE João Paulo da Cunha Faria \par}
      {\LARGE Maria Eduarda Teixeira Souza \par}
      \vfill

      {\large Divinópolis - MG \par}
      {\large Agosto de 2026 \par}
\end{titlepage}


\section{Introdução}

A representação digital de imagens é fundamentada na manipulação de cores através de espaços de cor bem definidos. O espaço RGB (Red, Green, Blue) é amplamente utilizado em computação gráfica por sua adequação a dispositivos de exibição como monitores e projetores, que emitem luz através da combinação de três canais de cor primária.

Este trabalho propõe a renderização da bandeira brasileira utilizando a biblioteca Pillow (\textit{Python Imaging Library}), explorando como primitivas gráficas simples e representações de cores em RGB podem ser combinadas para reproduzir elementos iconográficos bem conhecidos. A bandeira é composta por três camadas geométricas: um fundo verde que representa a história e a esperança, um losango amarelo que simboliza as riquezas minerais, e uma esfera azul (representação simplificada do globo terrestre) no centro.

\section{Fundamentos de Cor e Representação RGB}

\subsection{Espaço de Cor RGB}

O espaço RGB é um modelo aditivo de cores onde cada cor é representada como uma combinação de três componentes:
\begin{itemize}
    \item \textbf{Red (Vermelho):} intensidade de 0 a 255
    \item \textbf{Green (Verde):} intensidade de 0 a 255
    \item \textbf{Blue (Azul):} intensidade de 0 a 255
\end{itemize}

Uma tupla $(R, G, B)$ representa uma cor específica. Por exemplo:
\begin{itemize}
    \item Verde: $(0, 148, 64)$
    \item Amarelo: $(255, 255, 0)$
    \item Azul: $(0, 0, 255)$
\end{itemize}

Cada componente pode ser visto como um canal de luz que, quando combinado com os outros, produz a cor final perceida pelo observador. A biblioteca Pillow utiliza este padrão nativamente para objetos de imagem no modo \texttt{"RGB"}.

\subsection{Primitivas Gráficas}

O projeto utiliza duas primitivas fundamentais:

\begin{itemize}
    \item \textbf{Polígono:} Um polígono é definido por uma sequência de pontos no plano 2D. No contexto deste projeto, um losango é construído através de quatro pontos que definem seus vértices.
    \item \textbf{Elipse:} Uma elipse é definida por seu retângulo envolvente, especificado por dois pontos diagonais (superior esquerdo e inferior direito). Um círculo é um caso especial de elipse onde os eixos maior e menor possuem o mesmo comprimento.
\end{itemize}

\section{Metodologia}

\subsection{Materiais e Configuração}

O projeto foi desenvolvido em \textit{Python}, utilizando a biblioteca \texttt{Pillow} (PIL - Python Imaging Library), instalada via:

\begin{lstlisting}[style=pythonstyle]
pip install Pillow
\end{lstlisting}

A biblioteca Pillow fornece abstrações de alto nível para manipulação de imagens raster, permitindo criar, desenhar e salvar imagens em diversos formatos.

\subsection{Código Completo}

A seguir, apresenta-se a implementação completa do programa de renderização da bandeira brasileira:

\begin{lstlisting}[style=pythonstyle, caption={Código completo para renderização da bandeira brasileira}, label={lst:codigo_completo}]
from PIL import Image, ImageDraw

def criar_fundo_rgb(largura, altura, cor_rgb):
    """
    Cria uma imagem RGB com uma cor sólida
    """
    # Imagem.new aloca a matriz na memória.
    fundo = Image.new("RGB", (largura, altura), cor_rgb)

    # Retorna a imagem
    return fundo

def criar_losango_rgb(largura, altura, cor_rgb, bandeira):
    """
    Cria um losango RGB
    """

    # Ativando modo de sobreescrever pixels
    losango = ImageDraw.Draw(bandeira)

    # Coordenadas do losango
    pontoA = (largura / 2, altura * 0.05)
    pontoB = (largura * 0.95, altura / 2)
    pontoC = (largura / 2, altura * 0.95)
    pontoD = (altura * 0.05, altura / 2)

    # Lista de coordenadas para desenhar o losango
    coordenadas = [(pontoA), (pontoB), (pontoC), (pontoD)]

    # Criar o losango
    losango.polygon(coordenadas, fill=cor_rgb)

    # Retornar
    return bandeira


def criar_esfera_rgb(largura, altura, cor_rgb, bandeira):
    """
    Cria uma esfera RGB
    """

    # Ativando modo de sobreescrever pixels
    esfera = ImageDraw.Draw(bandeira)

    # Coordenadas da esfera (dois ponto, superior esquerdo e inferior direito)
    centro_x = largura / 2
    centro_y = altura / 2

    # Raio da esfera
    raio = largura * 0.20

    # Coordenadas do quadrado que envolve a esfera
    pontoA = (centro_x - raio, centro_y - raio)
    pontoB = (centro_x + raio, centro_y + raio)

    # Lista de coordenadas para desenhar a esfera
    coordenadas = [(pontoA), (pontoB)]

    # Criar a esfera
    esfera.ellipse(coordenadas, fill=cor_rgb)

    # Retornar
    return bandeira


if __name__ == "__main__":
    # dimensões da imagem
    LARGURA = 600
    ALTURA = 400

    # Testando a cor Verde: (0, 148, 64) RGB
    cor_verde_rgb = (0, 148, 64)
    print(f"Cor Verde em RGB: {cor_verde_rgb}")

    # Criar o fundo da bandeira
    bandeira = criar_fundo_rgb(LARGURA, ALTURA, cor_verde_rgb)

    # Testando a cor Amarela: (255, 255, 0) RGB
    cor_amarela_rgb = (255, 255, 0)
    print(f"Cor Amarela em RGB: {cor_amarela_rgb}")

    # Criar o losango
    bandeira = criar_losango_rgb(LARGURA, ALTURA, cor_amarela_rgb, bandeira)

    # Testando a cor Azul: (0, 0, 255) RGB
    cor_azul_rgb = (0, 0, 255)
    print(f"Cor Azul em RGB: {cor_azul_rgb}")

    # Criar a esfera
    bandeira = criar_esfera_rgb(LARGURA, ALTURA, cor_azul_rgb, bandeira)

    # Salvar a imagem
    bandeira.save("bandeira.png")
\end{lstlisting}

\subsection{Análise Detalhada do Código}

\subsubsection{Importações e Bibliotecas}

O programa inicia com a importação da biblioteca Pillow:

\begin{lstlisting}[style=pythonstyle]
from PIL import Image, ImageDraw
\end{lstlisting}

Dois componentes são importados:
\begin{itemize}
    \item \texttt{Image}: fornece a classe responsável por criar, carregar e manipular imagens raster
    \item \texttt{ImageDraw}: fornece ferramentas para desenhar primitivas gráficas (linhas, polígonos, elipses, etc.) sobre imagens existentes
\end{itemize}

\subsubsection{Função \texttt{criar\_fundo\_rgb}}

Esta função é responsável pela criação do fundo da bandeira:

\begin{lstlisting}[style=pythonstyle]
def criar_fundo_rgb(largura, altura, cor_rgb):
    fundo = Image.new("RGB", (largura, altura), cor_rgb)
    return fundo
\end{lstlisting}

Detalhamento:
\begin{itemize}
    \item \textbf{Parâmetros:} recebe largura, altura e uma tupla RGB
    \item \textbf{Image.new():} aloca uma nova imagem em memória no modo \texttt{"RGB"} com as dimensões especificadas
    \item \textbf{cor\_rgb:} define a cor inicial de todos os pixels da imagem
    \item \textbf{Retorno:} devolve o objeto \texttt{Image} criado para posterior manipulação
\end{itemize}

No contexto deste projeto, esta função é chamada com \texttt{cor\_verde\_rgb = (0, 148, 64)}, criando uma imagem $600 \times 400$ pixels totalmente preenchida com o verde da bandeira brasileira.

\subsubsection{Função \texttt{criar\_losango\_rgb}}

Esta função desenha o losango amarelo sobre a bandeira:

\begin{lstlisting}[style=pythonstyle]
def criar_losango_rgb(largura, altura, cor_rgb, bandeira):
    losango = ImageDraw.Draw(bandeira)
    pontoA = (largura / 2, altura * 0.05)
    pontoB = (largura * 0.95, altura / 2)
    pontoC = (largura / 2, altura * 0.95)
    pontoD = (altura * 0.05, altura / 2)
    coordenadas = [(pontoA), (pontoB), (pontoC), (pontoD)]
    losango.polygon(coordenadas, fill=cor_rgb)
    return bandeira
\end{lstlisting}

Detalhamento linha a linha:

\begin{itemize}
    \item \texttt{ImageDraw.Draw(bandeira):} cria um contexto de desenho vinculado à imagem passada como argumento. Este objeto permite executar operações gráficas sobre a imagem

    \item \texttt{pontoA = (largura / 2, altura * 0.05):} define o vértice superior do losango, posicionado no topo centro da imagem (50\% de largura, 5\% de altura)

    \item \texttt{pontoB = (largura * 0.95, altura / 2):} define o vértice direito do losango, posicionado no lado direito (95\% de largura, 50\% de altura)

    \item \texttt{pontoC = (largura / 2, altura * 0.95):} define o vértice inferior do losango, posicionado no fundo centro (50\% de largura, 95\% de altura)

    \item \texttt{pontoD = (altura * 0.05, altura / 2):} define o vértice esquerdo do losango (5\% de largura, 50\% de altura). Nota: há um erro nesta linha — deveria ser \texttt{(largura * 0.05)} ao invés de \texttt{(altura * 0.05)}, mas o código foi preservado como está na implementação original

    \item \texttt{coordenadas = [...]}: agrupa os quatro vértices em uma lista para passar ao método de desenho

    \item \texttt{losango.polygon(coordenadas, fill=cor\_rgb):} preenche o polígono definido pelos vértices com a cor especificada (amarelo)

    \item \texttt{return bandeira:} devolve a imagem modificada para posterior processamento
\end{itemize}

\subsubsection{Função \texttt{criar\_esfera\_rgb}}

Esta função desenha o círculo azul centralizado:

\begin{lstlisting}[style=pythonstyle]
def criar_esfera_rgb(largura, altura, cor_rgb, bandeira):
    esfera = ImageDraw.Draw(bandeira)
    centro_x = largura / 2
    centro_y = altura / 2
    raio = largura * 0.20
    pontoA = (centro_x - raio, centro_y - raio)
    pontoB = (centro_x + raio, centro_y + raio)
    coordenadas = [(pontoA), (pontoB)]
    esfera.ellipse(coordenadas, fill=cor_rgb)
    return bandeira
\end{lstlisting}

Detalhamento linha a linha:

\begin{itemize}
    \item \texttt{esfera = ImageDraw.Draw(bandeira):} cria um novo contexto de desenho sobre a imagem

    \item \texttt{centro\_x = largura / 2:} calcula a coordenada X do centro da imagem (300 pixels para largura 600)

    \item \texttt{centro\_y = altura / 2:} calcula a coordenada Y do centro da imagem (200 pixels para altura 400)

    \item \texttt{raio = largura * 0.20:} define o raio do círculo como 20\% da largura (120 pixels)

    \item \texttt{pontoA = (centro\_x - raio, centro\_y - raio):} calcula o canto superior esquerdo do retângulo envolvente do círculo (180, 80)

    \item \texttt{pontoB = (centro\_x + raio, centro\_y + raio):} calcula o canto inferior direito do retângulo envolvente (420, 320)

    \item \texttt{coordenadas = [(pontoA), (pontoB)]:} agrupa os dois pontos diagonais

    \item \texttt{esfera.ellipse(coordenadas, fill=cor\_rgb):} desenha uma elipse preenchida. Como os dois pontos diagonais possuem a mesma diferença em X e Y, cria-se um círculo perfeito preenchido com azul

    \item \texttt{return bandeira:} devolve a imagem modificada
\end{itemize}

\subsubsection{Bloco Principal}

O bloco \texttt{if \_\_name\_\_ == "\_\_main\_\_":} executa o fluxo de renderização quando o script é executado diretamente:

\begin{lstlisting}[style=pythonstyle]
LARGURA = 600
ALTURA = 400
cor_verde_rgb = (0, 148, 64)
print(f"Cor Verde em RGB: {cor_verde_rgb}")
bandeira = criar_fundo_rgb(LARGURA, ALTURA, cor_verde_rgb)
\end{lstlisting}

Define as dimensões da tela (600×400 pixels), especifica a cor verde e cria a imagem de fundo:

\begin{lstlisting}[style=pythonstyle]
cor_amarela_rgb = (255, 255, 0)
print(f"Cor Amarela em RGB: {cor_amarela_rgb}")
bandeira = criar_losango_rgb(LARGURA, ALTURA, cor_amarela_rgb, bandeira)
\end{lstlisting}

Define a cor amarela e desenha o losango sobre o fundo já criado:

\begin{lstlisting}[style=pythonstyle]
cor_azul_rgb = (0, 0, 255)
print(f"Cor Azul em RGB: {cor_azul_rgb}")
bandeira = criar_esfera_rgb(LARGURA, ALTURA, cor_azul_rgb, bandeira)
\end{lstlisting}

Define a cor azul e desenha o círculo sobre a bandeira:

\begin{lstlisting}[style=pythonstyle]
bandeira.save("bandeira.png")
\end{lstlisting}

Finalmente, salva a imagem resultante em formato PNG no diretório de trabalho atual.

\section{Resultados e Análises}

Durante a execução do programa principal, o sistema executa as seguintes operações em sequência:

\begin{enumerate}
    \item Cria uma imagem de $600 \times 400$ pixels com fundo verde RGB $(0, 148, 64)$
    \item Desenha um losango amarelo RGB $(255, 255, 0)$ sobre o fundo verde
    \item Desenha uma esfera azul RGB $(0, 0, 255)$ no centro da composição
    \item Salva a imagem final em formato PNG
\end{enumerate}

A saída da console exibe as cores em representação RGB:

\begin{verbatim}
    Cor Verde em RGB: (0, 148, 64)
    Cor Amarela em RGB: (255, 255, 0)
    Cor Azul em RGB: (0, 0, 255)
\end{verbatim}

A renderização visual resultante pode ser conferida na Figura \ref{fig:bandeira_final}, que apresenta a bandeira brasileira renderizada através de primitivas gráficas básicas. A composição em camadas demonstra como primitivas simples (polígono e elipse) combinadas com a manipulação precisa de cores RGB podem reproduzir artefatos visuais reconhecíveis.

\begin{figure}[H]
    \centering
    \includegraphics[width=0.9\textwidth]{images/bandeira.png}
    \caption{Renderização final da bandeira brasileira, composta por um fundo verde, um losango amarelo e uma esfera azul centralizada.}
    \label{fig:bandeira_final}
\end{figure}

\section{Conclusão}

O projeto demonstrou com sucesso a utilização da biblioteca Pillow para renderização de imagens através de primitivas gráficas e manipulação direta de cores no espaço RGB. A implementação confirmou que conceitos fundamentais de geometria computacional — posicionamento de vértices, cálculo de retângulos envolventes e preenchimento de formas — são suficientes para a composição de imagens reconhecíveis.

A abordagem modular adotada permite que cada camada da bandeira seja construída independentemente, facilitando ajustes de cores, dimensões e geometrias. Este trabalho reforça a importância da abstração entre modelos matemáticos de primitivas gráficas e sua implementação em bibliotecas de alto nível como Pillow.

\begin{thebibliography}{9}

\bibitem{rodrigo_primitivas}
RODRIGO, Raulivan. \textbf{Desenhar primitivas: Computação Gráfica}. Notas de Aula (05 - Desenhar primitivas.pdf). Departamento de Computação, Centro Federal de Educação Tecnológica de Minas Gerais (CEFET-MG), 2026.

\bibitem{pillow_docs}
CLARK, Alex; CONTRIBUTORS. \textbf{Pillow (PIL Fork) Documentation}. Disponível em: \url{https://pillow.readthedocs.io/}. Acesso em: Agosto de 2026.

\end{thebibliography}

\end{document}
